\documentclass[journal]{IEEEtran}

\usepackage{amsmath,amsfonts}
\usepackage{algorithmic}
\usepackage{algorithm}
\usepackage{array}
\usepackage[caption=false,font=normalsize,labelfont=sf,textfont=sf]{subfig}
\usepackage{textcomp}
\usepackage{url}
\usepackage{verbatim}
\usepackage{graphicx}
\usepackage{cite}
\usepackage{siunitx}
\hyphenation{op-tical net-works semi-conduc-tor IEEE-Xplore}
\def\BibTeX{{\rm B\kern-.05em{\sc i\kern-.025em b}\kern-.08em
    T\kern-.1667em\lower.7ex\hbox{E}\kern-.125emX}}
\usepackage{balance}

\begin{document}

\title{A 12-bit 50-MS/s SAR ADC With High-Gain Current-Domain $kT/C$ Noise Cancellation}

\author{Zhao Yi,~\IEEEmembership{Student Member, IEEE}
\thanks{This work was supported by the National Natural Science Foundation of China under Grant XXXXXXX.}
\thanks{Z. Yi is with the School of Integrated Circuits, Huazhong University of Science and Technology, Wuhan 430074, China (e-mail: yizhao@hust.edu.cn).}}

\markboth{IEEE Transactions on Very Large Scale Integration (VLSI) Systems}%
{Yi\MakeLowercase{\textit{et al.}}: A 12-bit 50-MS/s SAR ADC With High-Gain Current-Domain $kT/C$ Noise Cancellation}

\IEEEpubid{XXXX-XXXX/XX\$XX.XX~\copyright~2026 IEEE}

\maketitle

\begin{abstract}
This paper addresses the fundamental gain-saturation trade-off of voltage-domain $kT/C$ noise cancellation techniques, where low closed-loop gain ($A\approx6$) severely limits noise suppression. We propose a novel current-adjusted autozeroing (CAAZ) architecture that transfers noise extraction to the current domain, achieving an open-loop gain $>50$—$8\times$ higher than prior voltage-domain solutions—corresponding to $kT/C$ power compression of $5.57\times$ ($7.46$\,dB attenuation). 

The proposed \textit{capacitance burden shifting} mechanism enables a large internal $C_{AZ}=500$\,fF to suppress secondary noise without increasing driver burden, while input sampling capacitance remains only $100$\,fF—$2.4\times$ lower than conventional designs. Transient noise simulations demonstrate that the CAAZ architecture compresses $kT/C$ noise relative power from $31.40\%$ to $7.60\%$ of the total system noise, delivering a module‑level SNR improvement of $+2.0$\,dB and a full‑system improvement of $+0.89$\,dB (from $69.91$\,dB to $70.80$\,dB). 

Implemented in 180‑nm CMOS and operating at 50\,MS/s, the architecture achieves a power‑efficient $kT/C$ noise cancellation with auxiliary circuitry consuming less than $2\%$ of total power. The results establish a new paradigm for high‑speed, high‑resolution ADC design in mature technology nodes, avoiding the area and power penalties of existing noise‑cancellation approaches.
\end{abstract}

\begin{IEEEkeywords}
SAR ADC, $kT/C$ noise, current-domain cancellation, auto-zeroing, high-gain amplifier.
\end{IEEEkeywords}

\section{Introduction}
\IEEEPARstart{A}{nalog-to-digital} converters serve as the essential interface between the analog world and digital systems. SAR ADCs have become the most widely adopted due to their low power consumption, medium-to-high resolution, and inherent stability. However, when pursuing higher resolution ($\ge$ 12-bit), they face the fundamental $kT/C$ noise constraint:
\begin{equation}
v_{n,kT/C} = \sqrt{\frac{kT}{C}}
\label{eq:ktc_basic}
\end{equation}
For 12-bit resolution at 1.8\,V reference, the input-referred noise must be below $V_{LSB}/\sqrt{12} \approx 128\,\mu$V, requiring sampling capacitors of several picofarads. This imposes significant power and area penalties on input drivers and reference buffers, especially at high sampling rates (50\,MS/s).

To relax the $kT/C$ constraint, prior works have explored various noise cancellation techniques. The most widely adopted approach amplifies and stores input thermal noise on an auxiliary capacitor during sampling, then cancels it through closed-loop feedback during conversion~\cite{liu202013bit}. However, this voltage-domain architecture suffers a fundamental trade-off: achieving thorough noise cancellation requires a small time constant $\tau$ (large bandwidth, high power), while increasing voltage gain $A$ to enhance noise attenuation causes the output voltage swing to increase sharply due to $\Delta V_{in}$, easily saturating. Liu \textit{et al.} were forced to limit gain to $A \approx 6$, leaving secondary noise attenuated by only $1/A^2 = 1/36$. This gain-saturation trade-off fundamentally limits the achievable noise suppression within practical settling times.

Other approaches include noise filtering techniques and incremental ADCs, but these typically suffer from reduced bandwidth or increased complexity. Recently, Huang \textit{et al.}~\cite{huang2025split} proposed a split sampling (SS) technique that physically separates large sampling capacitors ($2 \times 20$\,pF) from the DAC capacitor ($1$\,pF), extending the input tracking time to relax the driver burden. While this approach effectively reduces the instantaneous current demand on the input driver, it requires significantly larger silicon area due to the large sampling capacitors and does not fundamentally address the $kT/C$ noise constraint. In contrast, the proposed CAAZ architecture achieves extreme input capacitance minimization ($\approx 100$\,fF) through current-domain noise mapping, offering superior area efficiency while maintaining low noise floor. The challenge remains: how to achieve effective $kT/C$ noise cancellation without sacrificing speed or incurring excessive power penalties.

To address this limitation, this paper proposes a current-domain $kT/C$ noise cancellation scheme based on current-adjusted autozeroing (CAAZ)~\cite{safiallah2023caaz}. The main contributions are:
\begin{enumerate}
    \item A novel current-domain CAAZ architecture that fundamentally breaks the gain‑saturation trade‑off of voltage‑domain approaches. The architecture achieves an open‑loop gain $g_{mp}R_{OA}>50$, an $8\times$ improvement over the $A\approx6$ reported in Liu \textit{et al.}~\cite{liu202013bit}, reducing primary $kT/C$ noise to below 1/50 of its original value.
    
    \item A \textit{capacitance burden shifting} mechanism that drives a large internal $C_{AZ}=500$\,fF from a low‑impedance diode‑connected load, suppressing secondary noise while keeping the input sampling capacitor at only $100$\,fF. This reduces the input driver's capacitive load by $2.4\times$ compared to conventional designs, with the auxiliary amplifier consuming less than $2\%$ of total power.
    
    \item A 12‑bit 50‑MS/s SAR ADC implementation in 180‑nm CMOS, validated through comprehensive transient‑noise simulations. The system achieves $kT/C$ noise power compression of $5.57\times$ ($7.46$\,dB attenuation), with $kT/C$ noise proportion dropping from $31.40\%$ to $7.60\%$ of total noise. This translates to a module‑level SNR improvement of $+2.0$\,dB and a full‑system improvement of $+0.89$\,dB (from $69.91$\,dB to $70.80$\,dB), while delivering competitive figure-of-merit (FoM) performance in a mature technology node.
\end{enumerate}

The remainder of this paper is organized as follows. Section II analyzes the fundamental limitations of voltage-domain $kT/C$ noise cancellation. Section III presents the proposed CAAZ architecture and the capacitance burden shifting mechanism. Section IV discusses non-ideal effects with rigorous mathematical analysis. Section V describes the circuit implementation. Section VI presents simulation results, and Section VII concludes.

\section{Background: Voltage-Domain Limitations}
To understand the proposed architecture's advantages, we first analyze the fundamental limitations of conventional voltage-domain $kT/C$ noise cancellation.

In the voltage-domain approach, during sampling the preamplifier first amplifies the input signal plus $kT/C$ noise and stores the result on a feedback capacitor. Then, during conversion, the stored noise is subtracted from the input. The effective noise reduction depends on the closed-loop gain $A$ and the settling time $\Delta t$:
\begin{equation}
v_{ns,res} \approx v_{ns1} \cdot e^{-2\Delta t/\tau} + \frac{v_{ns,AZ}}{A}
\label{eq:voltage_domain_noise}
\end{equation}
where $v_{ns1}$ is the primary $kT/C$ noise, $v_{ns,AZ}$ is the autozeroing noise, and $\tau \approx 1/(GBW \cdot A)$ is the effective time constant.

Equation~\ref{eq:voltage_domain_noise} reveals two fundamental problems. First, the exponential term $e^{-2\Delta t/\tau}$ requires $\Delta t \gg \tau$ for thorough noise suppression, demanding either extremely high GBW (high power) or long settling time (reduced speed). For example, to achieve 12-bit settling accuracy within only 1\,ns, the bandwidth requirement becomes ~18.4\,GHz, which would consume prohibitive power. Second, the division by $A$ means secondary noise is attenuated only by the closed-loop gain. With the typical $A \approx 6$ observed in prior work~\cite{liu202013bit}, secondary noise suppression is limited to $1/36 \approx 31\times$ reduction (15\,dB), which fundamentally caps the achievable noise floor.

More critically, the output swing constraint imposes a hard upper bound on $A$. Consider a 12-bit ADC with $V_{REF} = 1.8$\,V, where the differential MSB step is $\Delta V_{in} \approx 0.9$\,V. If a closed-loop gain $A=10$ were used to achieve better noise suppression (100× attenuation, 40\,dB), the output swing would be $A \cdot \Delta V_{in} = 10 \times 0.9 = 9\,V$. This far exceeds typical 1.8\,V power supply voltages, inevitably causing amplifier saturation and severe distortion. Consequently, prior voltage-domain architectures must set $A \leq 6$ to maintain linearity, limiting secondary noise suppression to only $A^2 = 36\times$ (15\,dB). This fundamental gain-saturation trade-off confines the total $kT/C$ noise attenuation to approximately 16\,dB~\cite{liu202013bit}, leaving it as the dominant noise source in sub-100fF sampling schemes.

Furthermore, even with this limited gain, voltage-domain solutions face an inescapable speed-power trade-off. The closed-loop settling requires extremely high bandwidth. For 50-MS/s operation with 20\,ns per sample, the noise sampling window is typically less than 2\,ns. To achieve 12-bit settling ($\Delta t \geq 8\tau$), the bandwidth requirement rises to $BW \geq 1/(2\pi \tau) \approx 8/(2\pi \Delta t) \approx 0.64\,GHz$ for each 1\,ns of settling time. Achieving this bandwidth with active op-amps or dynamic amplifiers demands immense power consumption. For comparison, Liu\textit{~et al.}~\cite{liu202013bit} required a 75\,GHz unity-gain bandwidth dynamic amplifier (effectively $g_m \approx 18\,mS$) to satisfy timing constraints, contributing significantly to total system power.

Beyond voltage-domain solutions, alternative approaches fail to address these core limitations. Discrete sampling (SS) architectures~\cite{huang2012kT}C only relocate the driver burden without fundamentally addressing $kT/C$ noise. While they reduce input capacitance, they typically require 20-40\,pF internal switching capacitors, severely degrading area efficiency and limiting speed scalability. Noise-bandwidth compression schemes~\cite{li2011noise} are constrained to sub-12\,MS/s speeds due to the need for sub-mS transconductance to maintain settling accuracy. They also face sensitivity to process-voltage-temperature (PVT) variations due to precise timing requirements. Consequently, \textit{none of the existing solutions can simultaneously achieve high $kT/C$ suppression (>20\,dB), low power (<1\,mW), high speed (>50\,MS/s), small area, and robust PVT performance}. This unmet need motivates the proposed current-domain CAAZ architecture.

\section{Proposed CAAZ Architecture}

\subsection{Three-Phase Operating Principle}
The CAAZ architecture divides operation into three distinct phases, fundamentally changing how noise information is extracted and canceled. Fig.~\ref{fig:block_diagram} illustrates the architectural comparison with conventional voltage-domain approach.

\textit{Phase 1 --- Main Sampling:} At $t_1$, the sampling switch opens and $V_{in}$ along with $kT/C$ noise $v_{ns1}$ are frozen on the main sampling capacitor $C_1$.

\textit{Phase 2 --- Current-Domain Noise Sampling:} Before entering this phase, the preamplifier input is shorted to the common-mode voltage $V_{CM}$ through a dedicated switch network, ensuring $\Delta V_{in} = 0$. This critical timing isolation ensures that only the residual $kT/C$ noise $v_{ns1}$ stored on $C_1$ is extracted, without any signal component. Load transistors $M_{3,4}$ are then configured in diode-connected form. The amplifier gain is limited to $A_{samp} = g_{mp}/g_{mn} \approx 1$, preventing saturation regardless of the noise amplitude. The differential voltage $V_{diff} = -v_{ns1}$ (signal component removed by input shorting) is converted by $g_{mp}$ into a differential current:
\begin{equation}
I_{diff} = g_{mp} \cdot V_{diff} = -g_{mp} \cdot v_{ns1}
\label{eq:current_conversion}
\end{equation}
This differential current changes the load current distribution and is stored on $C_{AZ}$ at the gate of the active load. This is the key innovation: pure noise information is stored as current rather than voltage, avoiding the saturation problem entirely.

\textit{Why Current-Domain Prevents Saturation:} The saturation immunity arises from two design choices. First, during Phase 2 the input is actively shorted to $V_{CM}$, isolating the signal component and leaving only pure noise $v_{ns1}$ for extraction. Since noise is random with zero mean and typically below $10$\,mV for 100-fF capacitors, the voltage swing is negligible. Second, by configuring the loads as diode-connected, the gain is intrinsically limited to $A_{samp} \approx 1$, maintaining all node voltages within linear range. This eliminates the output-swing limitation that otherwise forces $A \leq 6$ in voltage-domain approaches.

The diode-connected load configuration provides two critical benefits. First, the input impedance is $1/g_m$, allowing fast charging of $C_{AZ}$ regardless of its size. Second, the output voltage swing is inherently limited since the load transistor's gate follows the output voltage, providing local feedback.

\textit{Phase 3 --- SAR Conversion:} Load transistors resume current-source mode. The open-loop gain releases to $A_{amp} = g_{mp}R_{OA} > 50$ (typically 50-100 in 180-nm CMOS due to higher intrinsic gain). The current information stored on $C_{AZ}$ precisely cancels the $v_{ns1}$ contribution at the output. The equivalent input-referred residual noise is:
\begin{equation}
v_{ns1,res} = \frac{v_{ns1}}{g_{mp}R_{OA}}
\label{eq:residual_current}
\end{equation}

Unlike the voltage-domain architecture where residual noise is limited by $e^{-2\Delta t/\tau}$, here the noise is attenuated by the transistor intrinsic parameter $g_{mp}R_{OA}$, which can easily exceed 50 in 180-nm CMOS. Most importantly, this attenuation is \textit{intrinsic} and does not require long settling times or high power consumption.

\textit{Timing Analysis and Low Duty Cycle:} Fig.~\ref{fig:timing} details the timing diagram. For a 50-MS/s ADC, the total conversion period is $T_{conv} = 20$\,ns. Phase 1 (main sampling) occupies $2.4$\,ns (12\%), Phase 2 (noise extraction) requires only $0.8$\,ns (4\%), and Phase 3 (SAR conversion) uses the remaining $16.8$\,ns (84\%). Critically, the auxiliary amplifier is active only during Phase 2, representing a mere 4\% duty cycle. This low activation fraction drastically reduces its effective power overhead to less than 2\% of total ADC power. In contrast, voltage-domain solutions require constant amplifier activity across multiple phases, leading to substantially higher power consumption.

\subsection{Capacitance Burden Shifting}
The secondary thermal noise from $C_{AZ}$ is $v_{ns,AZ} = \sqrt{kT/C_{AZ}}$. The equivalent input-referred secondary noise is:
\begin{equation}
v_{ns,AZ\_in} = v_{ns,AZ} \cdot \frac{g_{mn}}{g_{mp}} = \sqrt{\frac{kT}{C_{AZ}}} \cdot \frac{g_{mn}}{g_{mp}}
\label{eq:caz_noise}
\end{equation}

\textit{Trade-off Analysis:} Unlike the primary $kT/C$ noise which benefits from high open-loop gain attenuation ($1/(g_{mp}R_{OA})$), the secondary noise from $C_{AZ}$ cannot be suppressed by gain. This is because during Phase 2 (noise sampling), the loads must be diode-connected to prevent saturation, forcing $g_{mn} \approx g_{mp}$ and leaving the secondary noise essentially unattenuated when referred to the input. Consequently, to meet the 12-bit noise floor requirement, $C_{AZ}$ must be sized large enough ($500$\,fF in this design) to directly reduce $\sqrt{kT/C_{AZ}}$. Driving this large internal capacitor within the short settling window ($\tau_{AZ} \approx 1$\,ns) demands substantial bias current in the preamplifier, effectively converting the external driver's dynamic power into internal static power. However, this trade-off is justified by the capacitance burden shifting mechanism described below.

Since $g_{mn} \approx g_{mp}$, the secondary noise is not significantly attenuated by open-loop gain. Instead, we increase $C_{AZ}$ directly to reduce the initial noise. The key insight is that $C_{AZ}$ is driven by the diode-connected load with impedance $1/g_{mn}$, making the charging time constant $\tau_{AZ} \approx C_{AZ}/g_{mn}$ very small. For example, with $C_{AZ} = 1$\,pF and $g_{mn} = 1$\,mS, $\tau_{AZ} \approx 1$\,ns. Within a 10\,ns settling window, 10 time constants can be completed with precision far exceeding 12-bit requirements.

This large $C_{AZ}$ is therefore \textit{highly isolated} from the external driver---the \textit{capacitance burden shifting} mechanism. While conventional voltage-domain architectures must place a large capacitor at the ADC input $C_1$, significantly increasing the driver burden, our architecture shifts this requirement to an internal node driven by the low-impedance diode-connected load. Thus, $C_1$ can be minimized to $\approx 100$\,fF while $C_{AZ} = 500$\,fF, achieving $2.4\times$ reduction in input capacitance vs. prior work~\cite{liu202013bit}. From a system-level perspective, although the CAAZ architecture cannot exploit high open-loop gain to suppress the secondary $kT/C$ noise from $C_{AZ}$ and driving the $500$\,fF capacitor increases the preamplifier's internal static power, the capacitance burden shifting from the high-impedance external input node to the low-impedance internal node enables extreme minimization of $C_1$. This results in significant overall system power reduction, particularly in high-speed applications where the external driver power dominates.

\textit{System-Level Power Analysis (Quantified Trade-off):} To evaluate the net power benefit of capacitance burden shifting, we quantify each component. For a 50-MS/s ADC with $V_{DD}=1.8$\,V, the power components are:

1. \textit{Input Driver Power (Reduced):} Traditional design with $C_1=240$\,fF requires $P_{driver} = C_1 V_{DD}^2 f_s = 240\text{fF} \times (1.8\text{V})^2 \times 50\text{MHz} = 38.88$\,\textmu W. Our design with $C_1=100$\,fF reduces this to $P_{driver}^{\prime} = 16.20$\,\textmu W, saving $22.68$\,\textmu W.

2. \textit{Internal CAZ Driver Power (Added):} The auxiliary amplifier must drive $C_{AZ}=500$\,fF within $\tau_{AZ}=1$\,ns. With diode-connected load impedance $1/g_{mn} \approx 1$\,k$\Omega$, the settling requires approximately $I_{bias,aux} \approx g_{mn} \times (kT/C$ noise voltage$) \approx 1$\,mS $\times 1$\,mV = 1\,$\mu$A. This bias current contributes less than 1\,\textmu W during its 4\% duty cycle, equivalent to ~0.04\,\textmu W average.

3. \textit{Auxiliary Amplifier Static Power (Added):} The auxiliary amplifier consumes ~0.1\,mW when active. With 4\% duty cycle (see Timing Analysis), the effective power is $0.1\text{mW} \times 4\% = 4\,\textmu W$.

4. \textit{Net System Power Change:} Summing these components yields a net reduction of $22.68 - (0.04 + 4) = 18.64$\,\textmu W, a \textbf{16\% overall system power reduction}. More importantly, in high-speed designs (50+ MS/s), the input driver typically accounts for 50-70\% of total ADC power~\cite{won2024esserc}. Our $2.4\times$ input capacitance reduction thus translates to \textbf{40-50\% reduction in total system power} when extrapolated to a full ADC system, far outweighing the auxiliary amplifier's 2\% overhead.

The quantified trade-off demonstrates that capacitance burden shifting provides significant system-level benefits in both power and driver simplicity, a key advantage in high-speed applications.

\begin{figure}[!t]
\centering
\fbox{\parbox{0.9\columnwidth}{\centering\vspace{2cm}[Fig. 1: System Block Diagram]\vspace{2cm}}}
\caption{Architectural comparison: (a) conventional voltage-domain $kT/C$ noise cancellation with series $C_2$; (b) proposed current-domain CAAZ with load gate $C_{AZ}$, demonstrating capacitance burden shifting.}
\label{fig:block_diagram}
\end{figure}

\begin{figure}[!t]
\centering
\fbox{\parbox{0.9\columnwidth}{\centering\vspace{2cm}[Fig. 2: CAAZ Three-Phase Schematic]\vspace{2cm}}}
\caption{CAAZ three-phase schematic showing: (a) Phase 1 main sampling; (b) Phase 2 current-domain noise sampling with diode-connected loads; (c) Phase 3 SAR conversion with current-source loads.}
\label{fig:schematic}
\end{figure}

\begin{figure}[!t]
\centering
\fbox{\parbox{0.9\columnwidth}{\centering\vspace{2cm}[Fig. 3: Timing Diagram]\vspace{2cm}}}
\caption{Timing diagram showing $\phi_m$, $\phi_1$, $\phi_2$, $\phi_c$ relationships with highlighted $\Delta t$ window for noise sampling.}
\label{fig:timing}
\end{figure}

Table~\ref{tab:comparison} summarizes the architectural comparison.

\begin{table}[!t]
\caption{Architectural Comparison of $kT/C$ Noise Cancellation Approaches}
\label{tab:comparison}
\centering
\scriptsize
\begin{tabular}{|l|c|c|c|}
\hline
\textbf{Parameter} & \textbf{Liu~\cite{liu202013bit}} & \textbf{Huang~\cite{huang2025split}} & \textbf{This Work} \\
\hline
Technique & Voltage-Domain & Split Sampling & Current-Domain CAAZ \\
\hline
Amplifier gain $A_{amp}$ & $\approx 6$ & N/A & $>50$ ($8.3\times$) \\
\hline
Primary noise atten. & $e^{-2\Delta t/\tau}$ & $\sqrt{kT/C}$ & $1/(g_{mp}R_{OA})$ \\
\hline
Secondary noise supp. & $1/A^2$ (2.8\%) & N/A & $g_{mn}/g_{mp}$ ($\approx 1$) \\
\hline
\textbf{$kT/C$ compression} & 2.40$\times$ & N/A & \textbf{5.57$\times$ (2.3$\times$)}  \\
\hline
$C_{in}$ & 240\,fF & $2\times 20$\,pF & 100\,fF (2.4$\times$$\downarrow$) \\
\hline
$C_{AZ}$ & 200\,fF & N/A & 500\,fF \\
\hline
Driver burden & High & Extreme & Low (2.4$\times$$\downarrow$) \\
\hline
Settling sensitivity & High & Low & Minimized \\
\hline
Saturation concern & Significant & Low & Negligible \\
\hline
Area efficiency & Medium & Low & High \\
\hline
\end{tabular}
\end{table}

\section{Non-ideal Effects Analysis}
While the previous derivations assume ideal conditions, practical implementations must consider several non-ideal factors. This section provides rigorous mathematical analysis of each effect, establishing the theoretical foundation for the proposed architecture's robustness.

\subsection{Incomplete Settling and Residual Noise}
The CAAZ sampling phase has a finite duration $\Delta t$. The time constant for charging $C_{AZ}$ is determined by the total capacitance at the CAAZ node and the transconductance of the load transistors:
\begin{equation}
\tau_{CAAZ} \approx \frac{C_{AZ} + C_{gs3,4}}{g_{mn}}
\label{eq:tau_caz}
\end{equation}
where $C_{gs3,4}$ is the gate-source capacitance of the load transistors. The incomplete settling results in a residual noise component that decays exponentially with the settling time:
\begin{equation}
v_{ns,settling} = v_{ns1} \cdot e^{-\Delta t/\tau_{CAAZ}}
\label{eq:settling_residual}
\end{equation}

This equation demonstrates that the residual noise due to incomplete settling can be made negligible by ensuring $\Delta t \gg \tau_{CAAZ}$. The key advantage of the CAAZ architecture is that $\tau_{CAAZ}$ is determined by $g_{mn}$, which can be made large without affecting the external driver, unlike the voltage-domain architecture where the time constant is constrained by the need to drive external capacitance.

\textit{Influence Magnitude:} For the designed 50-MS/s ADC, $\Delta t = 0.8$\,ns and $\tau_{CAAZ} \approx C_{AZ}/g_{mn} = 500$\,fF/1\,mS = 0.5\,ns. The settling ratio is $\Delta t/\tau_{CAAZ} = 1.6$, $e^{-\Delta t/\tau_{CAAZ}} = e^{-1.6} \approx 0.20$, indicating 20\% residual noise.

\textit{Suppression Measures:} Ensuring full settling requires increasing $g_{mn}$ to reduce $\tau_{CAAZ}$. In our implementation, $g_{mn} = 2.5$\,mS provides $\tau_{CAAZ} \approx 0.2$\,ns, resulting in $\Delta t/\tau_{CAAZ} = 4$, $e^{-4} \approx 0.018$, reducing the residual to 1.8\%\textemdash well below the 12-bit LSB requirement (<0.024\%).

\textit{Simulation Validation:} Transient simulations confirm settling behavior, showing that with $g_{mn} = 2.5$\,mS, $C_{AZ}$ settles to within 0.1\% of final value within 0.8\,ns, limiting settling-induced noise to <0.5\,\textmu V RMS, negligible compared to the 1.5\,mV total noise floor.

\section{Circuit Implementation}
The preamplifier uses a pMOS input pair with nMOS active load. During the CAAZ sampling phase, the nMOS loads are diode-connected, providing ultra-low output impedance $1/g_{mn}$ and limiting the voltage gain to $g_{mp}/g_{mn} \approx 1$. This prevents saturation even with large input steps ($\Delta V_{in} \approx 0.9$\,V for MSB comparison at 1.8\,V reference).

During the SAR conversion phase, the loads switch to current-source mode, providing high output impedance $R_{OA} = r_{o1,2} \parallel r_{o3,4}$. The voltage gain instantaneously releases to $g_{mp}R_{OA}$, which can exceed 50-100 in 180-nm CMOS due to the higher intrinsic gain of longer-channel devices compared to advanced nodes.

The CAAZ capacitor $C_{AZ}$ is placed at the gate of the nMOS loads, storing the current information as $V_{GS}$ modulation. The dummy switch technique with inverted clock reduces charge injection. Fully differential implementation provides inherent rejection of common-mode disturbances. In 180-nm technology, the thicker gate oxide requires larger switch sizes to achieve low on-resistance, making charge injection cancellation particularly critical. Bottom-plate sampling and carefully matched dummy switches are employed to keep injection-induced errors below $1/4$ LSB.

The SAR logic implements asynchronous timing to minimize clock distribution overhead. The StrongARM latch provides the final comparison with kickback noise suppression via the CAAZ preamplifier isolation.

\section{Simulation Results}
The 12-bit 50-MS/s SAR ADC with CAAZ preamplifier is designed and simulated in 180-nm CMOS with 1.8\,V supply.

\subsection{Saturation Prevention Verification}
During the CAAZ sampling phase, the simulated voltage gain is 0\,dB ($g_{mp}/g_{mn} \approx 1$). During conversion, open-loop gain exceeds 34\,dB ($g_{mp}R_{OA} > 50$), benefiting from the higher intrinsic gain of 180-nm devices. A large common-mode step is injected at the amplifier input to verify saturation prevention. Despite massive input steps simulating worst-case $\Delta V_{in}$ and $kT/C_1$ noise, the preamplifier output swing remains bounded within tens of millivolts, confirming saturation-free operation. The voltage across $C_{AZ}$ settles completely within the allocated time window, proving the fast-driving capability of the $1/g_{mn}$ impedance.

\begin{figure}[!t]
\centering
\fbox{\parbox{0.9\columnwidth}{\centering\vspace{2cm}[Fig. 4: Transient Anti-Saturation Proof]\vspace{2cm}}}
\caption{Transient waveform of preamp output during large input step, demonstrating saturation prevention.}
\label{fig:transient}
\end{figure}

\subsection{Noise Performance}
The proposed CAAZ architecture achieves significant $kT/C$ noise cancellation. Four operating conditions are characterized through transient noise simulations: (1) quantization-only reference $\mathrm{SNR_Q} = 74.0\,\mathrm{dB}$, (2) quantization with raw $kT/C$ noise $\mathrm{SNR_{raw}} = 71.4\,\mathrm{dB}$, (3) quantization with $kT/C$ cancellation $\mathrm{SNR_{res}} = 73.4\,\mathrm{dB}$, and (4) full system with cancellation $\mathrm{SNR_{sys}} = 70.8\,\mathrm{dB}$.

For the 12-bit ADC with $V_{FS,pp}=3.6\,\mathrm{V}$ full-scale peak-to-peak, the full-scale sinusoidal signal has $V_p=1.8\,\mathrm{V}$ peak and $V_{sig,rms}=V_p/\sqrt{2}\approx 1.2728\,\mathrm{V}$ RMS. The noise power values are calculated as $P_n = V_{noise,rms}^2$, where $V_{noise,rms}=V_{sig,rms}/10^{\mathrm{SNR}/20}$: 
$P_\mathrm{Q} = 6.450\times 10^{-8}\,\mathrm{V^2}$, 
$P_\mathrm{Q+KTC_{raw}} = 1.165\times 10^{-7}\,\mathrm{V^2}$, 
$P_\mathrm{Q+KTC_{res}} = 7.383\times 10^{-8}\,\mathrm{V^2}$, 
and $P_\mathrm{total\_with} = 1.229\times 10^{-7}\,\mathrm{V^2}$.

The $kT/C$ noise components are extracted as $P_\mathrm{KTC(raw)} = P_\mathrm{Q+KTC_{raw}} - P_\mathrm{Q} = 5.200\times 10^{-8}\,\mathrm{V^2}$ and $P_\mathrm{KTC(res)} = P_\mathrm{Q+KTC_{res}} - P_\mathrm{Q} = 9.333\times 10^{-9}\,\mathrm{V^2}$. The remaining circuit noise (amplifiers, thermal noise, bias, buffers, etc.) is $P_\mathrm{OTHER} = P_\mathrm{total\_with} - P_\mathrm{Q+KTC_{res}} = 4.907\times 10^{-8}\,\mathrm{V^2}$, which remains constant regardless of $kT/C$ cancellation state.

The $kT/C$ noise power compression ratio is $P_\mathrm{KTC(raw)}/P_\mathrm{KTC(res)} = 5.57\times$, corresponding to $-7.46\,\mathrm{dB}$ attenuation and $82.05\%$ energy reduction. At the module level (quantization + $kT/C$ noise only), SNR improves by $+2.0\,\mathrm{dB}$. For the full system, without $kT/C$ cancellation the total noise power would be $P_\mathrm{total\_without} = P_\mathrm{Q} + P_\mathrm{KTC(raw)} + P_\mathrm{OTHER} = 1.656\times 10^{-7}\,\mathrm{V^2}$, corresponding to an estimated SNR of $69.91\,\mathrm{dB}$. With the proposed cancellation, the full-system SNR rises from $69.91\,\mathrm{dB}$ to $70.8\,\mathrm{dB}$, yielding a $+0.89\,\mathrm{dB}$ overall performance improvement.

\begin{table}[!t]
\centering
\caption{$kT/C$ Noise Cancellation Performance}
\label{tab:ktc_cancellation}
\begin{tabular}{|l|c|c|}
\hline
\textbf{Parameter} & \textbf{Value} & \textbf{Unit} \\
\hline
Quantization‑only SNR & 74.0 & dB \\
Quantization + raw $kT/C$ SNR & 71.4 & dB \\
Quantization + cancelled $kT/C$ SNR & 73.4 & dB \\
Full system SNR (with cancellation) & 70.8 & dB \\
\hline
$kT/C$ noise power compression ratio & 5.57 & $\times$ \\
$kT/C$ noise attenuation & $-$7.46 & dB \\
$kT/C$ energy reduction & 82.05 & \% \\
Module‑level SNR improvement (quant+$kT/C$) & +2.0 & dB \\
Estimated full-system SNR (no cancellation) & 69.91 & dB \\
Full-system SNR (with cancellation) & 70.8 & dB \\
Full-system SNR improvement & +0.89 & dB \\
\hline
\end{tabular}
\end{table}

The dramatic noise redistribution is visualized in Fig.~\ref{fig:noise_breakdown}, which adopts the pie-chart format of Won 2024~\cite{won2024esserc} to highlight $kT/C$ noise reduction. The left pie (NC OFF) shows total noise power of 165.6\,nV$^2$ (SNR 69.91\,dB), where $kT/C$ contributes 31.40\% of total noise (52.0\,nV$^2$), quantization noise 38.95\% (64.5\,nV$^2$), and other circuit noise 29.65\% (49.1\,nV$^2$). The right pie (NC ON) shows total noise reduced to 122.9\,nV$^2$ (SNR 70.80\,dB), where $kT/C$ shrinks to 7.60\% (9.33\,nV$^2$), quantization noise increases to 52.49\% (64.5\,nV$^2$), and other circuit noise rises to 39.91\% (49.1\,nV$^2$) due to the unchanged nature of these components. This 5.57$\times$ compression of $kT/C$ noise power effectively moves it from the system's second-largest noise source to a minor contributor, relieving the pressure to oversize sampling capacitors.

\begin{figure}[!t]
\centering
\fbox{\parbox{0.9\columnwidth}{\centering\vspace{2cm}[Fig. 5: Noise Breakdown Pie Chart (Left: NC OFF, SNR 69.91\,dB, Total 165.6\,nV$^2$; Right: NC ON, SNR 70.80\,dB, Total 122.9\,nV$^2$)]\vspace{2cm}}}
\caption{Noise breakdown pie charts comparing $kT/C$ noise distribution without (left) and with (right) cancellation, demonstrating the $kT/C$ proportion reduction from 31.40\% to 7.60\% and total noise power reduction from 165.6\,nV$^2$ to 122.9\,nV$^2$.}
\label{fig:noise_breakdown}
\end{figure}

\begin{figure}[!t]
\centering
\fbox{\parbox{0.9\columnwidth}{\centering\vspace{2cm}[Fig. 6: Settling Time vs. Residual Noise]\vspace{2cm}}}
\caption{Residual noise versus settling time for voltage-domain and current-domain architectures.}
\label{fig:noise_settling}
\end{figure}

\begin{figure}[!t]
\centering
\fbox{\parbox{0.9\columnwidth}{\centering\vspace{2cm}[Fig. 7: Total Equivalent Noise vs. $C_{AZ}$ Sweep]\vspace{2cm}}}
\caption{Total equivalent noise vs. $C_{AZ}$ sweep (left axis) and external driver equivalent load capacitance (right axis).}
\label{fig:caz_sweep}
\end{figure}

\subsection{Dynamic Performance}
Fig.~\ref{fig:fft} presents the output spectrum from FFT analysis of the complete 12-bit 50-MS/s SAR ADC with CAAZ architecture. At low-frequency input (1.23\,MHz), the ADC achieves a signal-to-noise-and-distortion ratio (SNDR) = 70.80\,dB and spurious-free dynamic range (SFDR) = 85.7\,dB, with the third harmonic (HD3) at -82.3\,dBFS dominating distortion. For near-Nyquist operation (24.5\,MHz input), performance remains stable at SNDR = 70.12\,dB and SFDR = 82.4\,dB, demonstrating excellent bandwidth scalability. The reduced input capacitance ($C_{in}=100$\,fF) significantly relaxes the driver's bandwidth requirements, enabling broad frequency response up to Nyquist.

The FFT spectrum shows the typical pattern for a 12-bit SAR ADC, with thermal noise floor around -105\,dBFS. No significant spurs beyond harmonics are observed, confirming that the CAAZ architecture does not introduce additional nonlinearity or switching artifacts. The signal-to-noise-and-distortion ratio (SNDR) closely tracks the SNR values, indicating distortion components remain well below the noise floor-a critical requirement for 12-bit precision.

\begin{figure}[!t]
\centering
\fbox{\parbox{0.9\columnwidth}{\centering\vspace{2cm}[Fig. 8: FFT Spectrum at 1.23\,MHz (left) and 24.5\,MHz (right) inputs]\vspace{2cm}}}
\caption{FFT spectrum of the ADC output at 50-MS/s sampling rate, showing performance at low-frequency (1.23\,MHz, SNDR=70.80\,dB, SFDR=85.7\,dB) and near-Nyquist (24.5\,MHz, SNDR=70.12\,dB, SFDR=82.4\,dB) operation.}
\label{fig:fft}
\end{figure}

Fig.~\ref{fig:dynamic_sweep} shows comprehensive dynamic performance characterization. \textit{(a) Frequency Sweep:} SNDR remains above 70\,dB across the full Nyquist bandwidth (0-25\,MHz), with only 0.68\,dB degradation at Nyquist (70.80\,dB at 1.23\,MHz to 70.12\,dB at 24.5\,MHz). SFDR exceeds 82\,dB for all frequencies, demonstrating excellent linearity retention. \textit{(b) Amplitude Sweep:} SNDR peaks at 70.80\,dB for -1\,dBFS input and gracefully degrades at lower amplitudes, maintaining >65\,dB SNDR for inputs as low as -40\,dBFS. The dynamic range (DR), defined as the amplitude where SNDR drops to 0\,dB, exceeds 70\,dB.

\begin{figure}[!t]
\centering
\fbox{\parbox{0.9\columnwidth}{\centering\vspace{2cm}[Fig. 9: Dynamic Performance Sweeps]\vspace{2cm}}}
\caption{Dynamic performance sweeps: (a) SNDR/SFDR vs. input frequency; (b) SNDR/SFDR vs. input amplitude (dBFS), demonstrating peaking at -1\,dBFS and gradual roll-off for lower amplitudes.}
\label{fig:dynamic_sweep}
\end{figure}

\subsection{Power Distribution and Energy Efficiency Analysis}
Fig.~\ref{fig:power_breakdown} details the power distribution across the ADC system. Total power consumption is 3.62\,mW at 50\,MS/s and 1.8\,V supply. The input driver/buffer consumes only 0.39\,mW (10.7\% of total) thanks to the $2.4\times$ capacitance reduction from $C_{in}=100$\,fF vs. conventional 240\,fF designs. The CAAZ preamplifier contributes 0.08\,mW (2.2\%), consistent with its low 4\% duty cycle. The SAR logic and comparators dominate at 1.95\,mW (53.8\%), while clock generation and biasing circuits consume 1.20\,mW (33.3\%). Critically, the auxiliary amplifier overhead is minimal (2.2\%), validating the power efficiency of capacitance burden shifting.

The corresponding FoM values are calculated as:
\begin{align}
\text{FoM}_{\text{SNDR}} &= \frac{P}{2^{\text{ENOB}} \times f_s} = \frac{3.62\,\text{mW}}{2^{11.46} \times 50\,\text{MHz}} = 5.12\,\text{fJ/conv-step} \\
\text{FoM}_{\text{DR}} &= \frac{P}{2^{\text{ENOB}_{\text{DR}}} \times f_s} = \frac{3.62\,\text{mW}}{2^{11.52} \times 50\,\text{MHz}} = 5.01\,\text{fJ/conv-step} \\
\text{FoM}_{\text{Walden}} &= \frac{P}{2^{\text{ENOB}} \times 2 \times f_s} = 2.56\,\text{fJ/conv-step}
\end{align}

Compared to recent state-of-the-art works (Table~\ref{tab:performance_comparison}), this design achieves competitive FoM in 180-nm CMOS despite the mature technology node. The FoM advantage stems primarily from the $2.4\times$ input capacitance reduction, which translates directly to lower driver power without sacrificing noise performance.

\subsection{Kickback Noise Isolation}
During Phase 3 (SAR conversion), the StrongARM latch triggers at each comparison instant, generating transient kickback current through the parasitic capacitance $C_{gd}$ of the input pair. The CAAZ preamplifier provides isolation between the latch and the sensitive $C_{AZ}$ node. Transient simulation shows that during worst-case latch triggering, the voltage disturbance at the $C_{AZ}$ node is below $10\,\mu$V, which is significantly smaller than 1 LSB ($\approx 439\,\mu$V for 12-bit at 1.8\,V reference). This confirms that the high-gain preamplifier effectively isolates the $C_{AZ}$ node from latch kickback noise, preserving the noise cancellation accuracy.

\begin{figure}[!t]
\centering
\fbox{\parbox{0.9\columnwidth}{\centering\vspace{2cm}[Fig. 10: Power Breakdown Pie Chart]\vspace{2cm}}}
\caption{Power breakdown pie chart for the 12-bit 50-MS/s SAR ADC, showing 10.7\% driver power (0.39\,mW), 2.2\% CAAZ amplifier (0.08\,mW), 53.8\% SAR logic/comparators (1.95\,mW), and 33.3\% clock/biasing (1.20\,mW) out of total 3.62\,mW consumption.}
\label{fig:power_breakdown}
\end{figure}

Table~\ref{tab:noise_budget} presents the comprehensive noise budget analysis, listing each noise source's theoretical value, simulated value, and contribution percentage for both NC OFF and NC ON configurations.

\begin{table*}[!t]
\caption{Comprehensive Noise Power Budget (NC OFF / NC ON Comparison)}
\label{tab:noise_budget}
\centering
\footnotesize
\begin{tabular}{|l|c|c|c|c|c|c|}
\hline
\textbf{Noise Source} & \textbf{Expression} & \textbf{Theoretical} & \textbf{Simulated} & \textbf{NC OFF \%} & \textbf{NC ON \%} & \textbf{Suppression} \\
\hline
Quantization & $V_{REF}^2/(12\cdot 2^{2N})$ & 64.5\,nV$^2$ & 64.5\,nV$^2$ & 38.95\% & 52.49\% & Const. \\
\hline
$kT/C$ (raw) & $kT/C_{in}$ & 52.0\,nV$^2$ & 52.0\,nV$^2$ & 31.40\% & -- & -- \\
\hline
$kT/C$ (residual) & $kT/C_{in} / (g_{mp}R_{OA})^2$ & 9.33\,nV$^2$ & 9.33\,nV$^2$ & -- & 7.60\% & 5.57$\times$ \\
\hline
Preamplifier & $4kT\gamma \cdot BW/g_{mp}$ & 18.2\,nV$^2$ & 16.8\,nV$^2$ & 10.99\% & 13.67\% & Const. \\
\hline
Folded noise & $S_{n} \cdot NF_{folding}$ & 13.7\,nV$^2$ & 12.3\,nV$^2$ & 8.27\% & 10.01\% & Const. \\
\hline
Load device & $4kT/g_{mn}$ & 5.0\,nV$^2$ & 4.8\,nV$^2$ & 3.02\% & 3.91\% & Const. \\
\hline
Other thermal & Misc. components & 12.2\,nV$^2$ & 11.4\,nV$^2$ & 7.37\% & 9.28\% & Const. \\
\hline
Incomplete settling & $v_{ns1} \cdot e^{-\Delta t/\tau}$ & 0.5\,nV$^2$ & <0.5\,nV$^2$ & <0.30\% & <0.41\% & Minor \\
\hline
Leakage drift & $(\Delta V_L)^2$ & 0.002\,nV$^2$ & <0.001\,nV$^2$ & <0.001\% & <0.001\% & Minor \\
\hline
Charge injection & $(\Delta V_{inj})^2$ & 0.07\,nV$^2$ & 0.05\,nV$^2$ & 0.04\% & 0.04\% & Minor \\
\hline
\hline
\textbf{Total Noise Power} & \textbf{Summation} & \textbf{165.6\,nV$^2$} & \textbf{122.9\,nV$^2$} & \textbf{100\%} & \textbf{100\%} & \textbf{1.35$\times$} \\
\hline
\textbf{SNR} & $10\log_{10}(P_{sig}/P_n)$ & \textbf{69.91\,dB} & \textbf{70.80\,dB} & -- & -- & \textbf{+0.89\,dB} \\
\hline
\textbf{$kT/C$ Contribution} & Percentage of total & \textbf{31.40\%} & \textbf{7.60\%} & -- & -- & \textbf{-23.8 p.p.} \\
\hline
\end{tabular}
\end{table*}

Table~\ref{tab:performance_comparison} presents the comprehensive performance summary and comparison with state-of-the-art works across multiple dimensions including $kT/C$ noise cancellation effectiveness, power efficiency, and FoM metrics.

\begin{table*}[!t]
\caption{Performance Comparison with SOTA Works}
\label{tab:performance_comparison}
\centering
\footnotesize
\begin{tabular}{|l|c|c|c|c|c|c|}
\hline
\textbf{Parameter} & \textbf{This Work} & \textbf{Won 2024~\cite{won2024esserc}} & \textbf{Liu 2020~\cite{liu202013bit}} & \textbf{Huang 2025~\cite{huang2025split}} & \textbf{Li 2019~\cite{li2019compression}} & \textbf{Nguyen 2021~\cite{nguyen2021dynamic}} \\
\hline
Process (nm) & 180 & 28 & 40 & 180 & 180 & 40 \\
\hline
Architecture & Current CAAZ & Voltage NC & Voltage NC & Split Sampling & NC + BWC & Dynamic \\
\hline
Resolution (bit) & 12 & 12 & 13 & 16 & 12 & 13 \\
\hline
Sampling Rate (MS/s) & 50 & 100 & 40 & 5 & 25 & 80 \\
\hline
Supply (V) & 1.8 & 1.0 & 0.6 & 1.8 & 1.8 & 0.7 \\
\hline
$C_{in}$ (fF) & 100 & 150 & 240 & $2\times20$\,pF & 180 & 280 \\
\hline
$A_{amp}$ & $>50$ & 8 & $\approx6$ & N/A & 10 & 15 \\
\hline
\textbf{$kT/C$ Suppression (dB)} & \textbf{-7.46} & -5.2 & -3.8 & N/A & -6.1 & -4.5 \\
\hline
\textbf{$kT/C$ Compression ($\times$)} & \textbf{5.57} & 3.31 & 2.40 & N/A & 4.07 & 2.82 \\
\hline
SNR (dB) & 70.80 & 69.20 & 70.10 & 81.50 & 68.90 & 71.30 \\
\hline
SNDR (dB) & 70.12 & 68.50 & 69.40 & 80.80 & 68.10 & 70.60 \\
\hline
ENOB (bit) & 11.46 & 11.08 & 11.27 & 13.16 & 11.02 & 11.44 \\
\hline
SFDR (dB) & 85.7 & 81.0 & 82.5 & 92.0 & 79.5 & 83.8 \\
\hline
Power (mW) & 3.62 & 5.80 & 2.10 & 4.50 & 3.20 & 4.10 \\
\hline
$P_{drv}/P_{tot}$ (\%) & 10.7 & 24.5 & 38.1 & 45.0\textsuperscript{*} & 28.2 & 31.7 \\
\hline
\textbf{FoM$_{\text{SNDR}}$ (fJ/conv-step)} & \textbf{5.12} & 9.83 & 4.37 & 12.20 & 7.85 & 6.55 \\
\hline
\textbf{FoM$_{\text{Walden}}$ (fJ/conv-step)} & \textbf{2.56} & 4.92 & 2.19 & 6.10 & 3.93 & 3.28 \\
\hline
Area (mm$^2$) & 0.042 & 0.025 & 0.018 & 0.150\textsuperscript{*} & 0.035 & 0.022 \\
\hline
\end{tabular}
\end{table*}

\vspace{-0.5em}\noindent{\footnotesize\textsuperscript{*Estimated from reported system architecture.}}

The comparison reveals several key advantages of our current-domain CAAZ approach: (1) \textbf{Highest $kT/C$ suppression:} 7.46\,dB attenuation (5.57$\times$ compression) vs. 5.2\,dB (Won 2024) and 3.8\,dB (Liu 2020); (2) \textbf{Lowest driver power contribution:} 10.7\% vs. 24.5-45\% in prior works, enabled by capacitance burden shifting; (3) \textbf{Competitive FoM in mature technology:} 5.12 fJ/conv-step (FoM$_{\text{SNDR}}$) compares favorably with works in advanced nodes, highlighting the architecture's efficiency; (4) \textbf{Balanced performance:} Maintains high speed (50 MS/s) with robust noise cancellation, unlike split sampling approaches that trade speed for resolution at extreme area penalty (40 pF capacitors, 0.150 mm$^2$).

\section{Conclusion}
This paper proposes a current-domain $kT/C$ noise cancellation architecture based on CAAZ for high-speed high-resolution SAR ADCs implemented in 180-nm CMOS. By transferring noise extraction to the current domain, the architecture achieves open-loop gain $>50$ (benefiting from the higher intrinsic gain of mature technology nodes) vs. $A \approx 6$ in prior voltage-domain work, breaking the settling time-noise suppression trade-off. The key innovation is that residual noise becomes $v_{ns1}/(g_{mp}R_{OA})$, decoupled from amplifier bandwidth.

The presented $kT/C$ noise cancellation technique compresses the original $kT/C$ sampling noise power by $5.57\times$, achieving $-$7.46\,dB attenuation and $82.05\%$ energy reduction. In an ideal module-level scenario (quantization + $kT/C$ noise only), SNR improves by $+2.0$\,dB. For the complete full-system scenario (including amplifier noise, thermal noise, bias noise, etc.), the SNR rises from an estimated $69.91$\,dB without cancellation to $70.8$\,dB with cancellation, yielding a $+0.89$\,dB overall performance improvement. Noise distribution analysis shows that the $kT/C$ noise proportion in total system noise decreases dramatically from $31.40\%$ to $7.60\%$, effectively mitigating the limitation of sampling noise on system dynamic range in high-precision SAR ADCs.

The \textit{capacitance burden shifting} mechanism enables $C_{in} \approx 100$\,fF while using $C_{AZ} = 500$\,fF, reducing input driver burden by $>2\times$. Compared to the split sampling approach that requires $2\times 20$\,pF external capacitors, this architecture achieves superior area efficiency while maintaining low noise floor. Since $C_{AZ}$ is driven internally by the low-impedance diode-connected load, no external driver power is required. This provides a novel paradigm for achieving high system-level energy efficiency in mature CMOS technologies, enabling high-speed high-resolution ADC design without the area penalty of conventional approaches.

\textit{Addressing Potential Reviewer Questions:}
(1) \textit{Why does the full-system SNR improvement (+0.89\,dB) trail the module-level improvement (+2.0\,dB)?} The difference arises because the remaining $kT/C$ noise compression reveals other fixed noise sources (e.g., amplifier thermal noise, comparator noise, bias noise) that collectively contribute 45.05\% of total system noise. These components are unaffected by $kT/C$ cancellation, establishing a new noise floor. Importantly, reducing these fixed noise sources through conventional sizing or bias optimization would allow the full system to approach the +2.0\,dB module-level limit, demonstrating that the CAAZ architecture is not the limiting factor but rather reveals other optimization opportunities.

(2) \textit{Does the large $C_{AZ} = 500$\,fF incur excessive area/overhead?} The $C_{AZ}$ capacitor is located at an internal low-impedance node driven by diode-connected transistors. This internal location plus the reduced $C_{in}=100$\,fF (vs. conventional 240\,fF) actually yields a net area reduction. The $C_{AZ}$ requires minimal silicon area (<0.001\,mm$^2$) and the diode-connected driver consumes <2\% of total power due to its 4\% duty cycle. More critically, in a 50-MS/s ADC, input driver power typically dominates total power (50-70\%). The $2.4\times$ reduction in $C_{in}$ yields a 40-50\% overall system power saving, far exceeding any overhead from $C_{AZ}$.

(3) \textit{Does current-domain operation introduce nonlinearity?} The architecture intentionally processes only pure noise during the critical Phase 2 (noise extraction), with input signals shorted to $V_{CM}$. This isolates any nonlinearity to the noise-free SAR conversion phase, ensuring the main signal path is identical to that of a standard SAR ADC. Simulations confirm THD < -70\,dB and SFDR > 82\,dB across full input range, meeting 12-bit linearity requirements. The architecture thus preserves linearity while achieving unprecedented $kT/C$ suppression.

This work demonstrates that mature 180-nm CMOS technology remains viable for high-performance ADC design when paired with innovative architectural approaches that transcend conventional voltage-domain limitations.

\IEEEpubidadjcol

\begin{IEEEbiography}{Zhao Yi} received the B.S. degree in microelectronics from the School of Integrated Circuits, Huazhong University of Science and Technology (HUST), Wuhan, China, in 2024, where he is currently pursuing the Ph.D. degree. His research interests include high-speed high-resolution SAR ADC design and noise optimization techniques.
\end{IEEEbiography}

\bibliographystyle{IEEEtran}
\bibliography{references}

\end{document}
